\sectionheading{5.36 Python Operator}

Program-এ Data বা Value-এর উপর Calculation, Comparison,
Logical Test ইত্যাদি কাজ সম্পন্ন করার জন্য
\textbf{Operator} ব্যবহৃত হয়।

%================================================
% Definition
%================================================

\booksubheading{Operator কী?}

\begin{definitionbox}

এক বা একাধিক Value বা Variable-এর উপর নির্দিষ্ট Operation
সম্পন্ন করতে ব্যবহৃত Symbol বা Keyword-কে
\textbf{Operator} বলে।

\end{definitionbox}

উদাহরণ—

\begin{codebox}
\begin{english}
\ttfamily
a = 10\\
b = 5\\
c = a + b
\end{english}
\end{codebox}

\begin{center}

\textenglish{\ttfamily a, b}
\(\longrightarrow\)
\textbf{Operand}
\qquad
\textenglish{\ttfamily +}
\(\longrightarrow\)
\textbf{Operator}
\qquad
\textenglish{\ttfamily a + b}
\(\longrightarrow\)
\textbf{Expression}

\end{center}

\begin{rememberbox}
\begin{center}
\textbf{Operand}
\(\longrightarrow\)
\textbf{Operator দ্বারা Operation}
\(\longrightarrow\)
\textbf{Result}
\end{center}
\end{rememberbox}

%================================================
% Categories
%================================================

\booksubheading{Python Operator-এর প্রধান প্রকার}

\begin{tabularx}
{\linewidth}
{
|>{\bfseries\RaggedRight\arraybackslash}p{0.31\linewidth}
|Y|
}

\hline
\rowcolor{softBlue}

\multicolumn{1}{|c|}{\textbf{Operator-এর ধরন}}
&
\multicolumn{1}{c|}{\textbf{প্রধান Operator}}
\\
\hline

Arithmetic
&
\textenglish{\ttfamily + \quad - \quad * \quad / \quad // \quad \% \quad **}
\\
\hline

Comparison / Relational
&
\textenglish{\ttfamily == \quad != \quad > \quad < \quad >= \quad <=}
\\
\hline

Assignment
&
\textenglish{\ttfamily = \quad += \quad -= \quad *= \quad /= \quad //= \quad \%= \quad **=}
\\
\hline

Logical
&
\textenglish{\ttfamily and \quad or \quad not}
\\
\hline

Identity
&
\textenglish{\ttfamily is \quad is not}
\\
\hline

Membership
&
\textenglish{\ttfamily in \quad not in}
\\
\hline

Bitwise
&
\textenglish{\ttfamily
\& \quad \textbar \quad \textasciicircum \quad
\textasciitilde \quad << \quad >>}
\\
\hline

\end{tabularx}

%================================================
% Arithmetic
%================================================

\booksubheading{১. Arithmetic Operator}

সংখ্যার উপর গাণিতিক Operation করার জন্য
Arithmetic Operator ব্যবহৃত হয়।

ধরা যাক,
\textenglish{\ttfamily a = 10}
এবং
\textenglish{\ttfamily b = 3}।

\begin{tabularx}
{\linewidth}
{
|>{\bfseries\centering\arraybackslash}p{0.13\linewidth}
|>{\RaggedRight\arraybackslash}p{0.24\linewidth}
|>{\RaggedRight\arraybackslash}p{0.29\linewidth}
|Y|
}

\hline
\rowcolor{softBlue}

\textbf{Operator}
&
\textbf{কাজ}
&
\textbf{Example}
&
\textbf{Result}
\\
\hline

\textenglish{\ttfamily +}
& যোগ
& \textenglish{\ttfamily a + b}
& \textenglish{\ttfamily 13}
\\
\hline

\textenglish{\ttfamily -}
& বিয়োগ
& \textenglish{\ttfamily a - b}
& \textenglish{\ttfamily 7}
\\
\hline

\textenglish{\ttfamily *}
& গুণ
& \textenglish{\ttfamily a * b}
& \textenglish{\ttfamily 30}
\\
\hline

\textenglish{\ttfamily /}
& Division
& \textenglish{\ttfamily a / b}
& \textenglish{\ttfamily 3.333...}
\\
\hline

\textenglish{\ttfamily //}
& Floor Division
& \textenglish{\ttfamily a // b}
& \textenglish{\ttfamily 3}
\\
\hline

\textenglish{\ttfamily \%}
& Remainder
& \textenglish{\ttfamily a \% b}
& \textenglish{\ttfamily 1}
\\
\hline

\textenglish{\ttfamily **}
& Power
& \textenglish{\ttfamily a ** b}
& \textenglish{\ttfamily 1000}
\\
\hline

\end{tabularx}

%================================================
% Arithmetic examples side-by-side
%================================================

\noindent
\begin{minipage}[t]{0.49\linewidth}

\begin{sideexamplebox}{\textenglish{\ttfamily /} ও \textenglish{\ttfamily //}}

\begin{sidecodebox}
\begin{english}
\ttfamily
print(7 / 2)\\
print(7 // 2)
\end{english}
\end{sidecodebox}

\textbf{Output:}

\begin{center}
\textenglish{\ttfamily 3.5}
\qquad
\textenglish{\ttfamily 3}
\end{center}

\textenglish{\ttfamily /}
\(\longrightarrow\)
সাধারণ Division

\medskip

\textenglish{\ttfamily //}
\(\longrightarrow\)
Floor Division

\medskip

Python 3-এ
\textenglish{\ttfamily 10 / 2}
-এর Result
\textenglish{\ttfamily 5.0}।

\end{sideexamplebox}

\end{minipage}
\hfill
\begin{minipage}[t]{0.49\linewidth}

\begin{sideexamplebox}{\textenglish{\ttfamily \%} ও \textenglish{\ttfamily **}}

\begin{sidecodebox}
\begin{english}
\ttfamily
print(10 \% 3)\\
print(2 ** 3)
\end{english}
\end{sidecodebox}

\textbf{Output:}

\begin{center}
\textenglish{\ttfamily 1}
\qquad
\textenglish{\ttfamily 8}
\end{center}

Even Number পরীক্ষার শর্ত—

\begin{sidecodebox}
\begin{english}
\ttfamily
n \% 2 == 0
\end{english}
\end{sidecodebox}

\textenglish{\ttfamily \%}
\(\longrightarrow\)
Remainder

\medskip

\textenglish{\ttfamily **}
\(\longrightarrow\)
Power

\end{sideexamplebox}

\end{minipage}

\begin{rememberbox}

\textbf{\color{deepOrange}MCQ Trap}

\begin{center}

\textenglish{\ttfamily \%}
\(\longrightarrow\)
Remainder
\qquad
\textenglish{\ttfamily **}
\(\longrightarrow\)
Power

\qquad

\textenglish{\ttfamily \textasciicircum}
\(\longrightarrow\)
Bitwise XOR

\end{center}

\end{rememberbox}

%================================================
% String Operators
%================================================

\begin{rememberbox}

\textbf{\color{deepOrange}Python-এর বিশেষ ব্যবহার}

String-এর সঙ্গেও
\textenglish{\ttfamily +}
ও
\textenglish{\ttfamily *}
Operator ব্যবহার করা যায়।

\begin{center}

\textenglish{\ttfamily "ICT" + "Book"}
\(\longrightarrow\)
\textenglish{\ttfamily "ICTBook"}

\qquad

\textenglish{\ttfamily "Hi" * 3}
\(\longrightarrow\)
\textenglish{\ttfamily "HiHiHi"}

\end{center}

\textenglish{\ttfamily +}
\(\longrightarrow\)
String Concatenation

\qquad

\textenglish{\ttfamily *}
\(\longrightarrow\)
String Repetition

\end{rememberbox}

%================================================
% Comparison
%================================================

\booksubheading{২. Comparison / Relational Operator}

দুটি Value-এর মধ্যে তুলনা করার জন্য Comparison Operator
ব্যবহৃত হয়। Comparison-এর Result সাধারণত
\textenglish{\ttfamily True}
অথবা
\textenglish{\ttfamily False}।

ধরা যাক,
\textenglish{\ttfamily a = 10},
\textenglish{\ttfamily b = 5}।

\begin{tabularx}
{\linewidth}
{
|>{\bfseries\centering\arraybackslash}p{0.14\linewidth}
|>{\RaggedRight\arraybackslash}p{0.31\linewidth}
|>{\RaggedRight\arraybackslash}p{0.27\linewidth}
|Y|
}

\hline
\rowcolor{softBlue}

\textbf{Operator}
&
\textbf{অর্থ}
&
\textbf{Example}
&
\textbf{Result}
\\
\hline

\textenglish{\ttfamily ==}
& সমান কি না
& \textenglish{\ttfamily a == b}
& False
\\
\hline

\textenglish{\ttfamily !=}
& অসমান কি না
& \textenglish{\ttfamily a != b}
& True
\\
\hline

\textenglish{\ttfamily >}
& বড় কি না
& \textenglish{\ttfamily a > b}
& True
\\
\hline

\textenglish{\ttfamily <}
& ছোট কি না
& \textenglish{\ttfamily a < b}
& False
\\
\hline

\textenglish{\ttfamily >=}
& বড় অথবা সমান
& \textenglish{\ttfamily a >= b}
& True
\\
\hline

\textenglish{\ttfamily <=}
& ছোট অথবা সমান
& \textenglish{\ttfamily a <= b}
& False
\\
\hline

\end{tabularx}

\begin{rememberbox}

\textbf{\color{deepOrange}খুব গুরুত্বপূর্ণ}

\begin{center}

\textenglish{\ttfamily x = 10}
\(\longrightarrow\)
\textbf{Assignment}

\qquad

\textenglish{\ttfamily x == 10}
\(\longrightarrow\)
\textbf{Equality Check}

\end{center}

অর্থাৎ
\textenglish{\ttfamily =}
এবং
\textenglish{\ttfamily ==}
এক নয়।

\end{rememberbox}

%================================================
% Chained Comparison
%================================================

\begin{examplebox}{Python Chained Comparison}

Python-এ একই Expression-এ একাধিক Comparison লেখা যায়।

\begin{codebox}
\begin{english}
\ttfamily
x = 15\\
print(10 < x < 20)
\end{english}
\end{codebox}

\textbf{Output:}
\textenglish{\ttfamily True}

এর অর্থ—

\begin{center}

\textenglish{\ttfamily x > 10}
\qquad এবং \qquad
\textenglish{\ttfamily x < 20}

\end{center}

\end{examplebox}

%================================================
% Assignment
%================================================

\booksubheading{৩. Assignment Operator}

Variable-এর সঙ্গে Value Assign করার জন্য
Assignment Operator ব্যবহার করা হয়।

\begin{codebox}
\begin{english}
\ttfamily
x = 10
\end{english}
\end{codebox}

Compound Assignment একই Variable-এর উপর Operation করে
নতুন Value Assign করার সংক্ষিপ্ত পদ্ধতি।

\begin{tabularx}
{\linewidth}
{
|>{\bfseries\centering\arraybackslash}p{0.15\linewidth}
|>{\RaggedRight\arraybackslash}p{0.34\linewidth}
|Y|
}

\hline
\rowcolor{softBlue}

\textbf{Operator}
&
\textbf{উদাহরণ}
&
\textbf{সমতুল্য}
\\
\hline

\textenglish{\ttfamily +=}
&
\textenglish{\ttfamily x += 5}
&
\textenglish{\ttfamily x = x + 5}
\\
\hline

\textenglish{\ttfamily -=}
&
\textenglish{\ttfamily x -= 5}
&
\textenglish{\ttfamily x = x - 5}
\\
\hline

\textenglish{\ttfamily *=}
&
\textenglish{\ttfamily x *= 5}
&
\textenglish{\ttfamily x = x * 5}
\\
\hline

\textenglish{\ttfamily /=}
&
\textenglish{\ttfamily x /= 5}
&
\textenglish{\ttfamily x = x / 5}
\\
\hline

\textenglish{\ttfamily //=}
&
\textenglish{\ttfamily x //= 5}
&
\textenglish{\ttfamily x = x // 5}
\\
\hline

\textenglish{\ttfamily \%=}
&
\textenglish{\ttfamily x \%= 5}
&
\textenglish{\ttfamily x = x \% 5}
\\
\hline

\textenglish{\ttfamily **=}
&
\textenglish{\ttfamily x **= 2}
&
\textenglish{\ttfamily x = x ** 2}
\\
\hline

\end{tabularx}

উদাহরণ—

\begin{codebox}
\begin{english}
\ttfamily
x = 10\\
x += 5\\
print(x)
\end{english}
\end{codebox}

\textbf{Output:}
\textenglish{\ttfamily 15}

%================================================
% Logical
%================================================

\booksubheading{৪. Logical Operator}

একাধিক Condition যুক্ত করা বা Logical Result পরিবর্তনের জন্য
Logical Operator ব্যবহৃত হয়।

\begin{tabularx}
{\linewidth}
{
|>{\bfseries\centering\arraybackslash}p{0.18\linewidth}
|Y|
}

\hline
\rowcolor{softBlue}

\textbf{Operator}
&
\textbf{কাজ}
\\
\hline

\textenglish{\ttfamily and}
&
উভয় Condition সত্য হলে Result সত্য।
\\
\hline

\textenglish{\ttfamily or}
&
অন্তত একটি Condition সত্য হলে Result সত্য।
\\
\hline

\textenglish{\ttfamily not}
&
Logical Result বিপরীত করে।
\\
\hline

\end{tabularx}

%================================================
% Logical Examples
%================================================

\noindent
\begin{minipage}[t]{0.32\linewidth}

\begin{sideexamplebox}{\textenglish{\ttfamily and}}

\begin{sidecodebox}
\begin{english}
\ttfamily
age = 20\\
print(age >= 18\\
\hspace*{3mm}and age <= 30)
\end{english}
\end{sidecodebox}

Output:
\textenglish{\ttfamily True}

\end{sideexamplebox}

\end{minipage}
\hfill
\begin{minipage}[t]{0.32\linewidth}

\begin{sideexamplebox}{\textenglish{\ttfamily or}}

\begin{sidecodebox}
\begin{english}
\ttfamily
x = 5\\
print(x < 0 or x > 3)
\end{english}
\end{sidecodebox}

Output:
\textenglish{\ttfamily True}

\end{sideexamplebox}

\end{minipage}
\hfill
\begin{minipage}[t]{0.32\linewidth}

\begin{sideexamplebox}{\textenglish{\ttfamily not}}

\begin{sidecodebox}
\begin{english}
\ttfamily
x = True\\
print(not x)
\end{english}
\end{sidecodebox}

Output:
\textenglish{\ttfamily False}

\end{sideexamplebox}

\end{minipage}

%================================================
% Truth Table
%================================================

\begin{tabularx}
{\linewidth}
{
|>{\centering\arraybackslash}Y
|>{\centering\arraybackslash}Y
|>{\centering\arraybackslash}Y
|>{\centering\arraybackslash}Y|
}

\hline
\rowcolor{softBlue}

\textbf{A}
&
\textbf{B}
&
\textenglish{\ttfamily A and B}
&
\textenglish{\ttfamily A or B}
\\
\hline

True & True & True & True
\\
\hline

True & False & False & True
\\
\hline

False & True & False & True
\\
\hline

False & False & False & False
\\
\hline

\end{tabularx}

\begin{center}

\textenglish{\ttfamily not True}
\(\longrightarrow\)
\textenglish{\ttfamily False}

\qquad

\textenglish{\ttfamily not False}
\(\longrightarrow\)
\textenglish{\ttfamily True}

\end{center}

%================================================
% Identity + Membership
%================================================

\booksubheading{৫. Identity ও ৬. Membership Operator}

\noindent
\begin{minipage}[t]{0.49\linewidth}

\begin{sideexamplebox}{Identity Operator}

দুটি Variable একই Object-কে নির্দেশ করছে কি না পরীক্ষা করে।

\begin{center}
\textenglish{\ttfamily is}
\qquad
\textenglish{\ttfamily is not}
\end{center}

\begin{sidecodebox}
\begin{english}
\ttfamily
a = [1, 2]\\
b = a\\
print(a is b)
\end{english}
\end{sidecodebox}

\textbf{Output:}
\textenglish{\ttfamily True}

\medskip

\textenglish{\ttfamily ==}
\(\longrightarrow\)
Value সমান কি না

\medskip

\textenglish{\ttfamily is}
\(\longrightarrow\)
একই Object কি না

\end{sideexamplebox}

\end{minipage}
\hfill
\begin{minipage}[t]{0.49\linewidth}

\begin{sideexamplebox}{Membership Operator}

কোনো Value Collection-এর মধ্যে আছে কি না পরীক্ষা করে।

\begin{center}
\textenglish{\ttfamily in}
\qquad
\textenglish{\ttfamily not in}
\end{center}

\begin{sidecodebox}
\begin{english}
\ttfamily
numbers = [10, 20, 30]\\
print(20 in numbers)\\
print(50 not in numbers)
\end{english}
\end{sidecodebox}

\textbf{Output:}
\textenglish{\ttfamily True, True}

\medskip

String-এর ক্ষেত্রেও—

\begin{sidecodebox}
\begin{english}
\ttfamily
print("Py" in "Python")
\end{english}
\end{sidecodebox}

Output:
\textenglish{\ttfamily True}

\end{sideexamplebox}

\end{minipage}

\begin{rememberbox}

\textbf{\color{deepOrange}Membership কোথায় ব্যবহার করা যায়?}

\begin{center}

\textenglish{\ttfamily in / not in}
\(\longrightarrow\)
String, List, Tuple, Set, Dictionary

\end{center}

Dictionary-তে
\textenglish{\ttfamily in}
সাধারণত \textbf{Key} পরীক্ষা করে।

উদাহরণ—

\begin{center}

\textenglish{\ttfamily "name" in student}
\(\longrightarrow\)
Key আছে কি না

\end{center}

\end{rememberbox}

\begin{rememberbox}

\textbf{\color{deepOrange}ভুল করো না}

\begin{center}

\textenglish{\ttfamily ==}
\(\neq\)
\textenglish{\ttfamily is}

\qquad

\textenglish{\ttfamily in}
\(\longrightarrow\)
আছে কি না

\qquad

\textenglish{\ttfamily not in}
\(\longrightarrow\)
নেই কি না

\end{center}

Value তুলনা করার জন্য সাধারণত
\textenglish{\ttfamily ==}
ব্যবহার করা হয়।

\end{rememberbox}

%================================================
% Bitwise
%================================================

\booksubheading{৭. Bitwise Operator}

Integer-এর Binary Bit-এর উপর Operation করার জন্য
Bitwise Operator ব্যবহৃত হয়।

\begin{tabularx}
{\linewidth}
{
|>{\bfseries\centering\arraybackslash}p{0.16\linewidth}
|Y
|>{\bfseries\centering\arraybackslash}p{0.16\linewidth}
|Y|
}

\hline
\rowcolor{softBlue}

\textbf{Operator}
&
\textbf{নাম}
&
\textbf{Operator}
&
\textbf{নাম}
\\
\hline

\textenglish{\ttfamily \&}
&
Bitwise AND
&
\textenglish{\ttfamily \textbar}
&
Bitwise OR
\\
\hline

\textenglish{\ttfamily \textasciicircum}
&
Bitwise XOR
&
\textenglish{\ttfamily \textasciitilde}
&
Bitwise NOT
\\
\hline

\textenglish{\ttfamily <<}
&
Left Shift
&
\textenglish{\ttfamily >>}
&
Right Shift
\\
\hline

\end{tabularx}

\noindent
\begin{minipage}[t]{0.49\linewidth}

\begin{sideexamplebox}{Bitwise AND}

\begin{sidecodebox}
\begin{english}
\ttfamily
a = 5\\
b = 3\\
print(a \& b)
\end{english}
\end{sidecodebox}

\textbf{Output:}
\textenglish{\ttfamily 1}

\end{sideexamplebox}

\end{minipage}
\hfill
\begin{minipage}[t]{0.49\linewidth}

\begin{sideexamplebox}{Binary ধারণা}

\begin{sidecodebox}
\begin{english}
\ttfamily
5 = 101\\
3 = 011\\
---------\\
    001
\end{english}
\end{sidecodebox}

Result:
\textenglish{\ttfamily 1}

\end{sideexamplebox}

\end{minipage}

\begin{rememberbox}

\textbf{\color{deepOrange}Shift Operator-এর সহজ ধারণা}

Non-negative Integer-এর ক্ষেত্রে—

\begin{center}

\textenglish{\ttfamily x << n}
\(\longrightarrow\)
\(x \times 2^n\)

\qquad

\textenglish{\ttfamily x >> n}
\(\longrightarrow\)
\(x // 2^n\)

\end{center}

উদাহরণ—

\begin{center}

\textenglish{\ttfamily 5 << 1}
\(\longrightarrow\)
\textenglish{\ttfamily 10}

\qquad

\textenglish{\ttfamily 8 >> 1}
\(\longrightarrow\)
\textenglish{\ttfamily 4}

\end{center}

\end{rememberbox}

%================================================
% Unary and Binary
%================================================

\booksubheading{Unary ও Binary Operator}

Operand-এর সংখ্যার ভিত্তিতেও Operator-কে Unary ও Binary
হিসেবে বোঝা যায়।

\begin{tabularx}
{\linewidth}
{
|>{\bfseries\RaggedRight\arraybackslash}p{0.25\linewidth}
|>{\centering\arraybackslash}p{0.20\linewidth}
|Y|
}

\hline
\rowcolor{softBlue}

\textbf{ধরন}
&
\textbf{Operand}
&
\textbf{উদাহরণ}
\\
\hline

Unary
&
১টি
&
\textenglish{\ttfamily -x,\quad not x,\quad \textasciitilde x}
\\
\hline

Binary
&
২টি
&
\textenglish{\ttfamily a + b,\quad a > b}
\\
\hline

\end{tabularx}

উদাহরণ—

\begin{codebox}
\begin{english}
\ttfamily
x = 5\\
print(-x)
\end{english}
\end{codebox}

\textbf{Output:}
\textenglish{\ttfamily -5}

%================================================
% Precedence
%================================================

\booksubheading{Operator Precedence}

একটি Expression-এ একাধিক Operator থাকলে কোন Operator আগে
কাজ করবে তা \textbf{Operator Precedence} দ্বারা নির্ধারিত হয়।

\noindent
\begin{minipage}[t]{0.49\linewidth}

\begin{sideexamplebox}{Parentheses ছাড়া}

\begin{sidecodebox}
\begin{english}
\ttfamily
result = 2 + 3 * 4
\end{english}
\end{sidecodebox}

\textbf{Result:}
\textenglish{\ttfamily 14}

\medskip

কারণ
\textenglish{\ttfamily *}
আগে কাজ করে।

\end{sideexamplebox}

\end{minipage}
\hfill
\begin{minipage}[t]{0.49\linewidth}

\begin{sideexamplebox}{Parentheses সহ}

\begin{sidecodebox}
\begin{english}
\ttfamily
result = (2 + 3) * 4
\end{english}
\end{sidecodebox}

\textbf{Result:}
\textenglish{\ttfamily 20}

\medskip

কারণ Parentheses-এর অংশ আগে Evaluate হয়।

\end{sideexamplebox}

\end{minipage}

%================================================
% Precedence Table
%================================================

\begin{tabularx}
{\linewidth}
{
|>{\bfseries\RaggedRight\arraybackslash}p{0.25\linewidth}
|Y|
}

\hline
\rowcolor{softBlue}

\multicolumn{1}{|c|}{\textbf{অগ্রাধিকার}}
&
\multicolumn{1}{c|}{\textbf{Operator}}
\\
\hline

সর্বোচ্চ
&
Grouping:
\textenglish{\ttfamily ( )}
\\
\hline

এরপর
&
Exponentiation:
\textenglish{\ttfamily **}
\\
\hline

এরপর
&
Unary:
\textenglish{\ttfamily +,\;-,\;\textasciitilde}
\\
\hline

এরপর
&
\textenglish{\ttfamily *,\;/,\;//,\;\%}
\\
\hline

এরপর
&
\textenglish{\ttfamily +,\;-}
\\
\hline

এরপর
&
\textenglish{\ttfamily <<,\;>>}
\\
\hline

এরপর
&
\textenglish{\ttfamily \&}
\\
\hline

এরপর
&
\textenglish{\ttfamily \textasciicircum}
\\
\hline

এরপর
&
\textenglish{\ttfamily \textbar}
\\
\hline

এরপর
&
Comparison,
\textenglish{\ttfamily in},
\textenglish{\ttfamily is}
\\
\hline

এরপর
&
\textenglish{\ttfamily not}
\\
\hline

এরপর
&
\textenglish{\ttfamily and}
\\
\hline

তুলনামূলক কম
&
\textenglish{\ttfamily or}
\\
\hline

\end{tabularx}

\begin{rememberbox}

\textbf{\color{deepOrange}Exam-এর জন্য মূল Precedence}

\begin{center}

\textenglish{\ttfamily ( )}
\(\longrightarrow\)
\textenglish{\ttfamily **}
\(\longrightarrow\)
\textenglish{\ttfamily * / // \%}
\(\longrightarrow\)
\textenglish{\ttfamily + -}

\medskip

\textbf{Comparison}
\(\longrightarrow\)
\textenglish{\ttfamily not}
\(\longrightarrow\)
\textenglish{\ttfamily and}
\(\longrightarrow\)
\textenglish{\ttfamily or}

\end{center}

\end{rememberbox}

%================================================
% Associativity
%================================================

\booksubheading{Associativity}

একই Precedence-এর একাধিক Operator থাকলে কোন দিক থেকে
Operation সম্পন্ন হবে তা \textbf{Associativity} দ্বারা বোঝানো হয়।

\noindent
\begin{minipage}[t]{0.49\linewidth}

\begin{sideexamplebox}{সাধারণত Left to Right}

\begin{sidecodebox}
\begin{english}
\ttfamily
20 - 5 - 3
\end{english}
\end{sidecodebox}

অর্থাৎ—

\begin{center}
\textenglish{\ttfamily (20 - 5) - 3 = 12}
\end{center}

\end{sideexamplebox}

\end{minipage}
\hfill
\begin{minipage}[t]{0.49\linewidth}

\begin{sideexamplebox}{\textenglish{\ttfamily **} Right to Left}

\begin{sidecodebox}
\begin{english}
\ttfamily
2 ** 3 ** 2
\end{english}
\end{sidecodebox}

অর্থাৎ—

\begin{center}

\textenglish{\ttfamily 2 ** (3 ** 2)}
\(\longrightarrow\)
\textenglish{\ttfamily 2 ** 9}
\(\longrightarrow\)
\textenglish{\ttfamily 512}

\end{center}

\end{sideexamplebox}

\end{minipage}

\begin{rememberbox}

\textbf{\color{deepOrange}MCQ Trap:}
Exponentiation
\textenglish{\ttfamily **}
হলো \textbf{Right Associative}।

\end{rememberbox}

%================================================
% Operator Summary
%================================================

\chFiveInput{topics/36_python_operators/python_extensions.tex}

\booksubheading{Operator এক নজরে}

\begin{tabularx}
{\linewidth}
{
|>{\bfseries\RaggedRight\arraybackslash}p{0.23\linewidth}
|>{\RaggedRight\arraybackslash}p{0.36\linewidth}
|Y|
}

\hline
\rowcolor{softBlue}

\textbf{Category}
&
\textbf{Operator}
&
\textbf{প্রধান কাজ}
\\
\hline

Arithmetic
&
\textenglish{\ttfamily + - * / // \% **}
&
Calculation
\\
\hline

Comparison
&
\textenglish{\ttfamily == != > < >= <=}
&
Value তুলনা
\\
\hline

Assignment
&
\textenglish{\ttfamily = += -= *= ...}
&
Value Assign
\\
\hline

Logical
&
\textenglish{\ttfamily and or not}
&
Condition যুক্ত করা
\\
\hline

Identity
&
\textenglish{\ttfamily is,\;is not}
&
একই Object কি না
\\
\hline

Membership
&
\textenglish{\ttfamily in,\;not in}
&
Collection-এ আছে কি না
\\
\hline

Bitwise
&
\textenglish{\ttfamily
\& \quad \textbar \quad \textasciicircum
\quad \textasciitilde \quad << \quad >>}
&
Binary Bit Operation
\\
\hline

\end{tabularx}

%================================================
% Combined Example
%================================================

\booksubheading{Combined Example}

\begin{examplebox}{একই Program-এ একাধিক Operator}

\begin{codebox}
\begin{english}
\ttfamily
a = 10\\
b = 3\\
print(a + b)\\
print(a > b)\\
print(a \% 2 == 0)
\end{english}
\end{codebox}

\textbf{Output:}

\begin{center}

\textenglish{\ttfamily 13}
\qquad
\textenglish{\ttfamily True}
\qquad
\textenglish{\ttfamily True}

\end{center}

এখানে—

\begin{center}

\textbf{Arithmetic}
\(\longrightarrow\)
\textenglish{\ttfamily +,\;\%}

\qquad

\textbf{Comparison}
\(\longrightarrow\)
\textenglish{\ttfamily >,\;==}

\end{center}

\end{examplebox}

%================================================
% Common Mistakes
%================================================

\booksubheading{সাধারণ ভুল}

\begin{rememberbox}

\begin{itemize}

\item
\textenglish{\ttfamily =}
ও
\textenglish{\ttfamily ==}
-কে একই মনে করা।

\item
\textenglish{\ttfamily /}
ও
\textenglish{\ttfamily //}
-এর পার্থক্য না জানা।

\item
\textenglish{\ttfamily \%}
-কে Percentage Operator মনে করা; Python-এ এটি Remainder/Modulus Operator।

\item Power-এর জন্য
\textenglish{\ttfamily **}
-এর পরিবর্তে
\textenglish{\ttfamily \textasciicircum}
ব্যবহার করা।

\item
\textenglish{\ttfamily and},
\textenglish{\ttfamily or},
\textenglish{\ttfamily not}
-এর কাজ গুলিয়ে ফেলা।

\item
\textenglish{\ttfamily ==}
ও
\textenglish{\ttfamily is}
-কে একই মনে করা।

\item
\textenglish{\ttfamily in}
ও Identity Operator-এর কাজ গুলিয়ে ফেলা।

\item Operator Precedence উপেক্ষা করা।

\item Parentheses-এর প্রভাব না বোঝা।

\item
\textenglish{\ttfamily **}
-এর Right Associativity ভুলে যাওয়া।

\end{itemize}

\end{rememberbox}

%================================================
% Final Reminder
%================================================

\begin{rememberbox}

\textbf{\color{deepOrange}মনে রাখবে}

\begin{tabularx}
{\linewidth}
{
>{\bfseries\RaggedRight\arraybackslash}p{0.22\linewidth}
Y
>{\bfseries\RaggedRight\arraybackslash}p{0.22\linewidth}
Y
}

Arithmetic
&
\textenglish{\ttfamily + - * / // \% **}
&
Comparison
&
\textenglish{\ttfamily == != > < >= <=}
\\[0.6mm]

Assignment
&
\textenglish{\ttfamily =}
&
Logical
&
\textenglish{\ttfamily and or not}
\\[0.6mm]

Identity
&
\textenglish{\ttfamily is,\;is not}
&
Membership
&
\textenglish{\ttfamily in,\;not in}
\\[0.6mm]

Bitwise
&
\textenglish{\ttfamily
\& \;\textbar\; \textasciicircum\;
\textasciitilde\; <<\; >>}
&
Power
&
\textenglish{\ttfamily **}
\\[0.6mm]

Division
&
\textenglish{\ttfamily /}
&
Floor Division
&
\textenglish{\ttfamily //}
\\[0.6mm]

Remainder
&
\textenglish{\ttfamily \%}
&
Equality Check
&
\textenglish{\ttfamily ==}
\\

\end{tabularx}

\medskip

\begin{center}

\textenglish{\ttfamily =}
\(\neq\)
\textenglish{\ttfamily ==}
\qquad
\textenglish{\ttfamily ==}
\(\neq\)
\textenglish{\ttfamily is}

\medskip

\textenglish{\ttfamily \%}
\(\longrightarrow\)
Remainder
\qquad
\textenglish{\ttfamily **}
\(\longrightarrow\)
Power

\medskip

\textenglish{\ttfamily +}
\(\longrightarrow\)
String Concatenation-ও করতে পারে

\qquad

\textenglish{\ttfamily *}
\(\longrightarrow\)
String Repetition-ও করতে পারে

\medskip

\textbf{Precedence:}

\textenglish{\ttfamily ( )}
\(\longrightarrow\)
\textenglish{\ttfamily **}
\(\longrightarrow\)
\textenglish{\ttfamily * / // \%}
\(\longrightarrow\)
\textenglish{\ttfamily + -}

\medskip

\textbf{Comparison}
\(\longrightarrow\)
\textenglish{\ttfamily not}
\(\longrightarrow\)
\textenglish{\ttfamily and}
\(\longrightarrow\)
\textenglish{\ttfamily or}

\end{center}

\end{rememberbox}

\small

%================================================
% 5.37 Conditional Logic Statement
%================================================

